% Tree of Knowledge v4 --- causal locality reframing --- rewritten following LoRA structural patterns
% Core claim: specialization depth, adapter rank, and storage tier are one dimension.

\documentclass{article}
\usepackage[preprint,nonatbib]{neurips_2024}
\usepackage[utf8]{inputenc}
\usepackage[T1]{fontenc}
\usepackage{lmodern}
\usepackage{hyperref}
\usepackage{url}
\usepackage{booktabs}
\usepackage{amsfonts}
\usepackage{amsmath}
\usepackage{nicefrac}
\usepackage{microtype}
\usepackage{graphicx}
\usepackage{xcolor}
\usepackage{enumitem}
\usepackage{tikz}
\usetikzlibrary{arrows.meta}
\usepackage{wrapfig}

\newcommand{\placeholder}[1]{\begin{center}\fbox{\parbox{0.85\textwidth}{\centering\textcolor{red}{\textbf{[FIGURE: #1]}}}}\end{center}}
\newcommand{\tok}{Tree of Knowledge}

\title{Tree of Knowledge: Differentiating Fluid Intelligence \\ from Modular Crystallized Knowledge in Language Models}

\author{
  Erik de Bruijn \\
  Independent Researcher \\
  \texttt{erik@erikdebruijn.nl}
}

\begin{document}
\maketitle

% ============================================================================
\begin{abstract}
We separate a language model's \emph{fluid intelligence} (reasoning, syntax)
from \emph{crystallized intelligence} (domain knowledge) using variable-rank
LoRA adapters on FFN layers. With 655K trainable parameters on Qwen3-1.7B, we
reduce perplexity from 21.20 to 19.83---adapters of $\sim$45\,KB loadable from
NVMe in microseconds.

Experts readily specialize along \emph{functional} axes (structure vs.\
content) but not \emph{domain} axes---a pattern independently confirmed across
the MoE literature as the Standing Committee
effect~\cite{standing2026, domaindriver2026}. Removing either expert degrades
all domains equally ($\sim$1.3 PPL): the experts are genuinely specialized and
individually indispensable, but not independently removable. We show that \emph{causal locality} (selective ablation damage) rather than
weight orthogonality is the correct metric for modularity---and that it is
also the prerequisite for compute and memory efficiency, since only causally
local experts can be skipped or offloaded without uniform quality loss.
A Hessian analysis reveals that uniform routing is a \emph{saddle point}:
the LM loss encodes selective routing, but hard argmax blocks the gradient.
Domain selectivity ($M_{ij}$ CV) plateaus at $\sim$0.16 for domain-taxonomic
partitions, but a teacher-driven niche curriculum achieves strong routing selectivity
(92\% of niche tokens routed correctly) without corresponding ablation
damage---confirming that level-2 (routing) and level-3 (causal locality)
are genuinely distinct achievements.
\end{abstract}

% ============================================================================
\section{Introduction}
\label{sec:intro}

Current language models embed all knowledge---from common syntax to rare medical terminology---in the same parameters that perform reasoning. Adding knowledge means adding parameters, and therefore inference cost. The standard alternatives do not escape this trade-off: retrieval-augmented generation adds context tokens, tool calls add generation tokens, and longer contexts scale quadratically in attention. Every path to more knowledge costs inference budget.

The \tok{} framework aims to weaken this coupling by placing domain knowledge in separately-loadable adapters whose cost is zero until activated. The dense core handles reasoning and is always resident in VRAM. Crystallized knowledge lives in lightweight, hot-loadable LoRA adapters. Adding specialists does not slow the core's inference.

Prior work shows that learned models reside on a low intrinsic
dimension~\cite{li2018intrinsic,aghajanyan2021intrinsic}. We observe that
specialization depth, adapter rank, and storage tier are not three independent
design decisions---they correlate strongly. We distinguish two regimes:

\begin{itemize}[nosep,leftmargin=1em]
\item \textbf{Lens-like} (low rank, e.g.\ rank-1 to rank-4): the adapter does not alter the computation itself but \emph{redirects information} through the existing weights. The original FFN does the heavy lifting; the LoRA adapter bends the light. At 4--16\,KB, these corrections are always cached and cost nothing to load.
\item \textbf{Crystal-like} (high rank, e.g.\ rank-32 to rank-64): the $A \times B$ matrix is large enough to encode its own features---knowledge the original FFN does not possess. It is no longer a correction but an independent information carrier. At 128--256\,KB, these specialists are loaded from flash on demand.
\end{itemize}

The rank is not a hyperparameter to tune---it is a \emph{measurement} of the
dimensionality of the update-subspace needed beyond the parent representation.
A low rank means residual corrections fit in a small subspace (assimilation in
Piaget's terms~\cite{piaget1952}); a high rank means the expert requires
structurally new capacity that the parent cannot provide (accommodation).

\input{figure1.tex}

Standard Mixture-of-Experts (MoE) architectures attempt this separation but
fail in practice: load-balancing losses enforce near-uniform routing, preventing
any hot/cold distinction, and copy-initialized experts remain virtually
identical. We dedicate Section~\ref{sec:why-moe-fails} to this failure mode,
which motivates every design choice in \tok{}.

\paragraph{Concrete benefits.}
\begin{itemize}
    \item Specialists measured in \textbf{kilobytes}, not megabytes (4\,KB--256\,KB vs.\ 16\,MB per MoE expert).
    \item Trainable on \textbf{commodity hardware}, in parallel---each adapter trains independently.
    \item \textbf{Composable by a community}---contributors add specialists without retraining the core.
    \item \textbf{No inference-latency overhead}---the core runs at full speed; adapters merge into the forward pass.
\end{itemize}

Throughout this paper, we use \textbf{Qwen3-1.7B} (28 layers, 2.0B parameters)
as our running example, training on C4 data. All reported numbers are from this
model unless otherwise noted.

% ============================================================================
\section{Why Can't Standard MoE Differentiate?}
\label{sec:why-moe-fails}

The promise of Mixture-of-Experts is that different experts will learn different
things. In practice, they do---but not in the way we expect. Experts readily
specialize along \emph{functional} axes (structure vs.\ content, punctuation
vs.\ rare words) but resist specializing along \emph{domain} axes (medical
vs.\ legal vs.\ code). The word ``expert'' evokes a human domain specialist,
but MoE experts behave more like specialized tools: one handles
fasteners, another handles cutting, regardless of whether the project is
carpentry or plumbing. We spent six experimental arms and 40+ hours of GPU
training on Qwen3-1.7B discovering why domain specialization is harder than
functional specialization, and what it takes to cross that boundary.

\paragraph{The copy-initialization trap.}
The standard upcycling recipe~\cite{komatsuzaki2023sparse} creates experts by
copying the dense model's FFN weights. Each expert starts in the same basin of
attraction. Our decisive experiment: we forced completely disjoint data
partitions across experts, training for 12M tokens. Result: expert weight cosine
similarity remained at \textbf{0.998}. The copy initialization creates a basin
that disjoint data alone cannot escape on practical timescales.

\paragraph{Load-balancing destroys structure.}
MoE training adds an auxiliary load-balancing loss that penalizes uneven routing.
This is necessary to prevent expert collapse, but it has a side effect: it
enforces near-uniform token distribution across experts, eliminating the very
frequency structure that would allow hot/cold tiering. Every expert sees the
same distribution and converges to the same function.

\paragraph{Six arms, one conclusion.}
Table~\ref{tab:lessons} summarizes six experimental arms that systematically
explored alternatives. Prescriptive cost losses create routing structure (Gini
0.60) but the effect dilutes over training. Descriptive losses fail entirely
(Gini 0.05). Biological mechanisms help marginally (Gini 0.46). Without
explicit divergence pressure on the \emph{weights themselves}, no genuine
specialization emerges.

\begin{table}[t]
\centering
\caption{Six experimental arms on Qwen3-1.7B that establish what does \emph{not}
produce genuine expert differentiation.}
\label{tab:lessons}
\begin{tabular}{@{}llp{6.5cm}@{}}
\toprule
Arm & Key metric & Lesson \\
\midrule
A: Prescriptive cost loss  & Gini 0.60$\to$0.49 & Creates routing skew, but skew dilutes over training \\
B: Reuse-distance proxy     & Gini 0.05           & Descriptive loss produces no structure \\
C: Adaptive + cost          & 25 rounds, Gini dilutes & Piaget mechanism works, but routing Gini is wrong metric \\
D: Pure data-driven         & Gini 0.11           & No specialization without resource constraint \\
E: Bio-integrated           & Gini 0.46           & Six bio mechanisms help modestly, not fundamentally \\
F: KD-Warm upcycle          & CosSim 0.998        & Copy-init experts don't differentiate via data alone \\
\bottomrule
\end{tabular}
\end{table}

These failures taught us that expert differentiation requires three forces acting
\emph{together}: (1)~different starting points (not copy-initialization),
(2)~different training data, and (3)~explicit divergence pressure on the weights.
The \tok{} framework instantiates all three through learntropy-driven progressive
splitting with contrastive LoRA adapters.

\paragraph{Three levels of differentiation.}
Our experiments reveal that ``expert specialization'' is not binary but admits
at least three distinct levels, each progressively harder to achieve:

\begin{enumerate}[nosep]
    \item \textbf{Parameter differentiation.} Expert weight matrices become
          orthogonal in weight space (low CosSim). Achievable via contrastive
          loss. \emph{Necessary but not sufficient.}
    \item \textbf{Routing differentiation.} The router sends different token
          populations to different experts based on meaningful features (not
          just load balancing). Requires informative hidden states at the
          routing point (Section~\ref{sec:divergence}).
    \item \textbf{Causal modularity.} Each expert must be selectively
          indispensable for a specific token population: removing it causes
          concentrated damage to that population while leaving others
          unaffected. This is the requirement for hot-pluggability, and it is
          fundamentally different from generic otherness---experts must not
          merely \emph{differ}, they must \emph{own} distinct subsets of the
          model's competence. Measured by the ablation matrix $M_{ij}$, not by
          weight geometry.
\end{enumerate}

Our framework achieves level~1 reliably (CosSim $<$0.3) but has not yet
achieved levels~2 or~3. The gap between these levels is the central open
problem this paper identifies.

\paragraph{Domains may not be the natural primitive.}
A recurring theme across our experiments and the literature~\cite{standing2026,
domaindriver2026, openmoe2024} is that functional axes (syntax, token type,
sentence position) are the \emph{primary} basis of expert specialization, while
domain axes (medical, legal, code) are secondary and may not emerge as a
natural partition at all. If this is correct, the right framing is not ``how
do we make experts specialize by domain'' but ``what are the natural factors
along which expert specialization occurs, and are domains a meaningful
downstream consequence of those factors?''

\paragraph{Unifying efficiency mechanisms in a single structure.}
Several efficiency mechanisms that MoE systems often introduce as separate
architectural components---such as shared always-on
experts~\cite{deepseek2024moe}, driver-like experts with disproportionate
causal influence~\cite{domaindriver2026}, and selectively routed
specialists---may instead correspond to different depths in a single tree. The
trunk provides a shared substrate; early, low-rank forks can serve as
lightweight infrastructure; and deeper, higher-rank branches can implement
conditional residual specialization. This remains a simplifying hypothesis
rather than a demonstrated result: our current experiments have not yet
produced the multi-level tree needed to test it directly.

The tree also resembles a broader pattern seen in biological systems: shared
early processing followed by progressively more specialized downstream
computation. We do not claim a mechanistic correspondence to cortical
organization, but the analogy is suggestive: both biological and artificial
systems appear to benefit from hierarchical specialization under resource
constraints~\cite{bena2025modularity}.

% ============================================================================
\section{Method: Progressive LoRA Forking}
\label{sec:method}

The \tok{} architecture has three components: a shared trunk, a branching
mechanism, and a divergence loss. We describe each briefly; mathematical details
are in Appendix~\ref{app:math}.

\subsection{LoRA on FFN Layers, Not Attention}

A key design choice: we apply LoRA adapters to the \emph{feed-forward network}
(specifically \texttt{gate\_proj} and \texttt{up\_proj}) rather than to
attention layers. While Hu et al.~\cite{hu2022lora} demonstrated LoRA
primarily on attention projections, mechanistic interpretability research
suggests that FFN layers function as key-value memories that store factual
associations~\cite{geva2021transformer}, whereas attention layers perform
routing and compositional binding. Our goal is to modularize knowledge---whether that modularization ultimately
follows domain boundaries (medical vs.\ legal) or functional boundaries
(structure vs.\ content) is an empirical question. The FFN is the natural
target because we are modifying what the model \emph{stores}, not how it
\emph{routes} information. The attention layers remain shared across all
experts, preserving general reasoning and compositional abilities (fluid
intelligence).

\subsection{Shared Trunk, Branching Depth}

The trunk boundary is determined empirically by layer divergence analysis
(Section~\ref{sec:divergence}). For Qwen3-1.7B (28 layers), domain-specific
signals emerge at layer~18 (CKA increases 8$\times$ between layers 14 and 18).
We partition the transformer into:
\begin{itemize}
    \item \textbf{Layers 0--17 (trunk):} Shared FFN, always VRAM-resident.
          Fluid intelligence core---general reasoning, syntax, attention
          patterns. $\sim$1.2B parameters.
    \item \textbf{Layers 18--27 (expert layers):} Variable-rank LoRA adapters
          on FFN. Each adapter $\sim$45\,KB at rank-4 per layer.
\end{itemize}
Total specialist parameters: 655K for 2 experts across 10 layers, versus
800+\,MB for full expert copies in standard MoE.

\subsection{Learntropy-Driven Mitosis}

% TODO: Learntropy definition is provisional. Wozniak (SuperMemo) independently
% coined "learntropy" as the attractiveness of a signal to the learn drive —
% a NET reward signal (reward minus decoding-failure penalty) that depends on
% prior knowledge and can go negative. Our current definition (raw CE loss) is
% a simplification: it does not distinguish "productively hard" (Wozniak's
% positive learntropy) from "impossibly hard" (negative learntropy). Bimodality
% detection partially recovers this distinction, but a proper learntropy
% function should model net learning value, not raw difficulty. We are working
% toward a definition closer to Wozniak's, operationalized as a loss function.
% See: https://supermemo.guru/wiki/Learntropy
%
\textbf{Definition: Learntropy.} We use \emph{learntropy} to denote the per-token prediction difficulty experienced by a specific expert:
\begin{equation}
\mathcal{L}_{\text{learn}}(e_i, t) = -\log P(t \mid \text{context}, e_i)
\end{equation}
where $e_i$ is expert $i$ and $t$ is the predicted token. The term was independently coined by Wozniak~\cite{wozniak2018learntropy} to denote the attractiveness of a signal to the learn drive system---a net reward that depends on prior knowledge and can go negative when material exceeds the learner's capacity. Our operationalization as raw cross-entropy loss is a simplification: it measures difficulty but does not distinguish productive difficulty (material in the learner's zone of proximal development~\cite{vygotsky1978}) from noise. Bimodality detection (below) partially recovers this distinction by identifying when an expert's loss distribution contains both learnable and unlearnable populations.

We use learntropy as a \emph{diagnostic signal for the learning system itself}: not only as a training objective but as an instrument that determines when structural change is needed. This reflects Shannon's insight that information content drives the allocation of representational resources~\cite{shannon1948}---and Piaget's observation that prediction failure (accommodation) triggers the creation of new cognitive structure~\cite{piaget1952}.

When an expert's learntropy distribution is bimodal---some tokens easy, others hard---the expert must either adapt or divide. Following biological cells, we distinguish two responses:

\textbf{Gene expression (rank increase).} The expert first attempts to handle both populations by increasing its rank---activating more parameters, like a cell expressing additional genes to handle a broader range of signals. This is the cheaper response: no routing overhead, no data partitioning.

\textbf{Mitosis (expert split).} If rank increase does not resolve the bimodality, the expert undergoes \emph{mitosis}: it divides into two daughter experts, each inheriting the parent's weights plus a small random perturbation. Like biological daughter cells, they start genetically identical but \emph{differentiate} through exposure to different environments (data partitions) and explicit pressure to diverge (contrastive loss). The procedure:
\begin{enumerate}
    \item Score each expert's tokens by learntropy.
    \item Test for bimodality (Hartigan's dip test).
    \item If bimodal and rank increase insufficient: duplicate the LoRA adapter.
          Each daughter inherits parent weights plus perturbation, starting at rank-1.
    \item Daughter experts \emph{differentiate} through their data environment and contrastive loss (Section~\ref{sec:contrastive}).
\end{enumerate}
The decision to divide is triggered by an information-theoretic signal, not a schedule---analogous to how cell division is triggered by growth signals, not a timer.

\subsection{Variable Rank as Accommodation}

The rank of an adapter measures the dimensionality of its update-subspace
relative to the parent representation. We formalize rank growth as Piagetian
\emph{accommodation}: structural expansion that occurs when new residual
structure cannot be \emph{assimilated} within the existing subspace.

For an expert $e$ with current update-subspace $U_e$ (spanned by its LoRA
matrices $A$ and $B$), let $g$ be the gradient on the residual error. Decompose
$g = g_\parallel + g_\perp$ where $g_\parallel \in U_e$ and $g_\perp \perp U_e$.
The \emph{accommodation ratio} is:
\begin{equation}
\mathcal{A}(e) = \frac{\|g_\perp\|}{\|g\|}
\label{eq:accommodation}
\end{equation}

When $\mathcal{A}(e)$ is low, the error fits within the current subspace
(assimilation---no structural change needed). When $\mathcal{A}(e)$ is high,
the error requires capacity the expert does not have. The response depends on
the \emph{structure} of $g_\perp$: if it is unimodal (one coherent residual
population), rank growth is appropriate; if it is multimodal (distinct
populations requiring different corrections), splitting is appropriate.

\textbf{Practical caveat.} At low rank, $\mathcal{A}(e)$ is dominated by
geometry: a rank-$r$ subspace in $d$ dimensions captures only $r/d$ of random
directions, yielding $\mathcal{A}(e) \approx 1 - r/d$ even for uninformative
gradients. We measured $\mathcal{A}(e) \approx 0.99$ at rank~4 across all
training phases---close to the random baseline of $1 - 4/2048 = 0.998$. A
useful trigger therefore requires either higher rank (so the subspace covers a
non-trivial fraction) or a \emph{normalized} accommodation ratio that
compares against the rank-appropriate baseline.

We distinguish two regimes:
\begin{itemize}
    \item \textbf{Rank 1--4} (4--16\,KB): A lens. Residual corrections fit
          in a small subspace; the parent representation is nearly sufficient.
    \item \textbf{Rank 16--64} (64--256\,KB): A crystal. The update-subspace
          has grown large enough to encode structure the parent cannot represent.
\end{itemize}

\begin{table}[t]
\centering
\caption{Specialization depth, adapter rank, and storage tier correlate.}
\label{tab:unification}
\begin{tabular}{@{}llll@{}}
\toprule
Specialization & Adapter & Size & Storage tier \\
\midrule
Minimal (lens) & Rank-1 LoRA & 4\,KB & Always resident \\
Moderate       & Rank-16 LoRA & 64\,KB & VRAM cache \\
Substantial (crystal) & Rank-64 LoRA & 256\,KB & NVMe (instant load) \\
Full expert    & Dense copy & 12\,MB & NVMe (prefetchable) \\
\bottomrule
\end{tabular}
\end{table}

\subsection{Contrastive Divergence Loss}
\label{sec:contrastive}

Without explicit pressure, siblings drift toward the shared optimum (we observed
CosSim 0.998). The contrastive loss provides the missing force:
\begin{equation}
\mathcal{L}_{\text{contrast}} = \lambda_c \sum_{i < j} \max\bigl(0,\; \cos(\theta_{ij}) - \tau\bigr)
\label{eq:contrastive}
\end{equation}
where $\theta_{ij}$ is the angle between sibling weight vectors and $\tau$ is a
target maximum similarity. Combined with disjoint data and LoRA initialization,
this produces the differentiation that eluded all six prior arms.

\subsection{Teacher-Guided Data Selection via Zone of Proximal Development}
\label{sec:zpd}

Contrastive loss differentiates expert \emph{weights}, but does not guarantee
that experts encode different \emph{knowledge}. Two experts can be orthogonal
in weight space yet both required for every input---as we observe with the
first-level Structure/Content split (Section~\ref{sec:hotloading}).

To produce modular experts whose removal causes \emph{selective} damage, each
expert must train on fundamentally different data. We select this data using a
teacher model and Vygotsky's Zone of Proximal Development (ZPD): the set of
examples that are learnable but not yet learned.

For each text chunk $x$ and expert $e_i$, we compute:
\begin{equation}
\text{ZPD}(x, e_i) = \frac{\text{PPL}_{\text{student}}(x, e_i)}{\text{PPL}_{\text{teacher}}(x)}
\label{eq:zpd}
\end{equation}
A high ZPD score means the student struggles where the teacher succeeds---the
knowledge exists in the teacher but is absent from the student. A low score
means either both struggle (noise, not knowledge) or both succeed (nothing to
learn). Each expert receives the chunks where \emph{its own} ZPD is highest,
producing genuinely disjoint training distributions.

We tested ZPD scoring with Qwen3-30B as teacher and Qwen3-1.7B as student on
500 C4 chunks: Spearman $\rho = 0.958$ between teacher and student difficulty
rankings, with only 14.2\% of chunks in the ZPD zone. In this configuration,
the teacher adds little signal beyond what the student already knows about its
own difficulty. This does not falsify the ZPD concept itself---the result is
specific to this model pair on generic C4 data. ZPD-guided selection may prove
more effective (a)~after experts have diverged enough to have genuinely
different difficulty profiles, (b)~with a much larger teacher, or (c)~on
domain-specific data where teacher--student knowledge gaps are wider.

% ============================================================================
\section{Experiments}
\label{sec:experiments}

All experiments use Qwen3-1.7B on C4 data. We build incrementally: each
subsection adds one capability and reports the result.

\subsection{Layer Divergence Analysis: Where to Fork}
\label{sec:divergence}

Before placing adapters, we measure where domain-specific representations
emerge in the unmodified model. We pass 200 samples from each of five C4
domains (medical, code, legal, conversational, news) through Qwen3-1.7B and
compute per-layer Centered Kernel Alignment (CKA) between domain-specific
hidden state distributions.

\begin{table}[h]
\centering
\caption{Domain divergence by layer (selected). CKA measures alignment between
domain-specific hidden state distributions (higher = more domain-specific).}
\label{tab:divergence}
\begin{tabular}{@{}lrrr@{}}
\toprule
Layer range & CKA & Fisher ratio & Interpretation \\
\midrule
0--8 (early) & 0.003--0.004 & 0.004--0.005 & Domain-agnostic \\
9--14        & 0.005--0.013 & 0.006--0.009 & Minimal divergence \\
15--17       & 0.020--0.062 & 0.010--0.012 & Transition zone \\
\textbf{18--21} & \textbf{0.106--0.200} & \textbf{0.014--0.016} & \textbf{Domain-specific} \\
22--27       & 0.206--0.216 & 0.015--0.018 & Strong divergence \\
\bottomrule
\end{tabular}
\end{table}

CKA increases by an order of magnitude between layers 14 and 18 (0.013 to
0.106). This empirically determines the trunk boundary: layers below 18 carry
insufficient domain signal for expert specialization. Our initial experiments
(below) used a trunk of 14 layers;
the layer divergence analysis motivated extending the trunk to 18 layers in
subsequent experiments.

\subsection{Baseline}

\begin{table}[h]
\centering
\caption{Qwen3-1.7B baseline on C4 validation.}
\label{tab:baseline}
\begin{tabular}{@{}lr@{}}
\toprule
Configuration & Perplexity \\
\midrule
Dense baseline (no adapters) & 21.20 \\
\bottomrule
\end{tabular}
\end{table}

The dense Qwen3-1.7B achieves a perplexity of 21.20 on the C4 validation set.
This is the number every subsequent experiment must match or improve.

\subsection{Phase 1: First Split with LoRA Adapters}

We freeze the trunk (layers 0--14) and attach rank-4 LoRA adapters to layers
15--27, training on full C4 data without domain partitioning. This validates
that the LoRA+trunk architecture can recover baseline quality.

\begin{table}[h]
\centering
\caption{Phase 1: LoRA adapters on the upper layers.}
\label{tab:phase1}
\begin{tabular}{@{}lrr@{}}
\toprule
Configuration & PPL & $\Delta$ from baseline \\
\midrule
Dense baseline & 21.20 & --- \\
Trunk + LoRA (layers 15--27) & 19.59 & $-$1.61 \\
\bottomrule
\end{tabular}
\end{table}

The LoRA adapters not only recover baseline quality but \emph{improve} it,
reducing perplexity to 19.59. The trunk preserves reasoning capacity while the
adapters specialize the upper layers.

\subsection{Phase 2: Contrastive Forking}

We create two sibling LoRA adapters at the first fork point, train them on
disjoint data partitions with the contrastive loss (Eq.~\ref{eq:contrastive}),
and measure differentiation.

\begin{table}[h]
\centering
\caption{Phase 2: Sibling differentiation with contrastive loss.}
\label{tab:phase2}
\begin{tabular}{@{}lrrr@{}}
\toprule
Configuration & CosSim & PPL & Note \\
\midrule
Disjoint data only (no contrast) & 0.975 & 19.62 & Minimal differentiation \\
Disjoint data + contrastive loss & 0.840 & 19.58 & Strong differentiation \\
\bottomrule
\end{tabular}
\end{table}

The contrastive loss reduces cosine similarity from 0.975 to 0.840---a
qualitative shift from ``nearly identical'' to ``measurably different
functions''---while \emph{improving} perplexity to 19.58. Differentiation does
not trade off against quality.

\subsection{Progressive Forking}
\label{sec:forking}

With the trunk boundary informed by the layer divergence analysis
(Section~\ref{sec:divergence}), we compare forking at two boundary points.

\begin{table}[h]
\centering
\caption{Trunk boundary comparison: Phase 2 results (2 experts, rank 4,
contrastive weight 0.1). Layer-18 trunk places experts in the domain-specific
zone identified by the divergence analysis.}
\label{tab:trunk-comparison}
\begin{tabular}{@{}lrrrr@{}}
\toprule
Configuration & Expert layers & CosSim & PPL & Params \\
\midrule
Trunk-14 (layers 15--27) & 14 & 0.840 & 19.58 & 1.9M \\
Trunk-14 + Phase 3 (rank 32) & 14 & 0.278 & 19.33 & 11.2M \\
Trunk-14 + Phase 3 (rank 4) & 14 & 0.433 & 19.45 & 1.9M \\
\textbf{Trunk-18 (layers 18--27)} & \textbf{10} & \textbf{0.771} & \textbf{20.14} & \textbf{655K} \\
Trunk-18 + Phase 3 (rank 4) & 10 & 0.311 & 19.83 & 655K \\
\bottomrule
\end{tabular}
\end{table}

The trunk-18 configuration achieves deeper differentiation (CosSim 0.311 at
rank 4) than trunk-14 at the same rank (0.433), with 3$\times$ fewer
parameters. Both fork location and rank constraint affect differentiation:
trunk-14 at rank 32 reaches 0.278, while trunk-14 at rank 4 reaches only
0.433. Trunk-18 at rank 4 matches the rank-32 result (0.276 vs.\ 0.278),
suggesting that \emph{placing experts in the domain-divergent zone compensates
for 8$\times$ fewer parameters}.

However, routing analysis (Section~\ref{sec:magnitude}) shows that none of
these configurations achieve domain-specific routing. Weight-space
differentiation (low CosSim) does not imply functional specialization. Recent
work on MoE expert roles~\cite{standing2026} confirms this pattern at scale:
a ``standing committee'' of experts handles functional/structural processing
across all domains, while domain-specific routing is peripheral.

No second fork (2$\to$4 experts) has triggered in any configuration on
generic C4 data.

\subsection{Hot-Loading Validation}
\label{sec:hotloading}

Our first hot-loading test \emph{failed}: removing either first-level expert
degraded all text types equally ($\sim$1.3 PPL each). Token routing analysis
revealed Expert~0 handles rare/content words (71\% of rare tokens) while
Expert~1 handles punctuation and function words (86\% of punctuation). Every
sentence requires both structure and content---neither expert is independently
removable.

\textbf{Observation:} The first split is functional (structure vs.\ content),
not domain-specific. Both experts are always needed.

\textbf{Hypothesis (untested):} Domain-level modularity may emerge at deeper
tree levels, where sub-experts within the Content branch specialize on
different domains. If so, removing a domain sub-expert would selectively
degrade that domain's PPL. This requires a second fork, which has not yet
occurred in our experiments.

\textbf{Alternative explanations we cannot yet rule out:}
\begin{itemize}
    \item Deeper splits may also produce functional rather than domain
          differentiation, in which case the tree never becomes modular.
    \item The structure/content split may be an artifact of contrastive loss
          pushing toward the easiest orthogonal direction in weight space,
          not a meaningful linguistic distinction.
    \item The tree metaphor may not accurately describe the learning dynamics
          at all---the branching may reflect optimization landscape geometry
          rather than knowledge organization.
\end{itemize}

\begin{table}[h]
\centering
\caption{Ablation matrix $M_{ij}$ for trunk-18 Phase~3 final checkpoint
(CosSim 0.311). Each cell shows $\Delta$PPL when expert $j$ is removed,
evaluated on domain $i$. A modular system would show high values on the
diagonal and low values off-diagonal.}
\label{tab:ablation-matrix}
\begin{tabular}{@{}lrr@{}}
\toprule
Domain $i$ & Remove E0 & Remove E1 \\
\midrule
Medical       & +1.38 & +1.16 \\
Code          & +1.63 & +1.47 \\
Legal         & +1.42 & +1.20 \\
News          & +1.26 & +1.05 \\
Conversational & +1.22 & +1.05 \\
\midrule
\textbf{Mean} & \textbf{+1.38} & \textbf{+1.19} \\
\bottomrule
\end{tabular}
\end{table}

The ablation matrix is uniform: coefficient of variation 0.11 across domains
(threshold for meaningful selectivity: $>$0.3). Each expert accounts for a substantial fraction of the total LoRA improvement
(high substitutability: removing either costs $\sim$1.2--1.4 PPL of the
$\sim$1.4 total benefit). This confirms that
weight orthogonality (CosSim 0.311) does not produce causal locality---the
experts are near-interchangeable general-purpose adapters, not domain
specialists.

We also tested a shared + gated-routed architecture inspired by
DeepSeekMoE~\cite{deepseek2024moe}: one always-active expert per layer plus
one gated expert with a learned sigmoid gate and sparsity penalty. The gate
converged to 0.985 (near-always-on) with domain selectivity of 0.007 (needed
$>$0.20). Given the option of using additional capacity, the model uses it
uniformly rather than selectively.

\subsection{Hessian Analysis of Routing Landscape}
\label{sec:hessian}

To determine whether uniform routing is a local minimum (basin) or a saddle
point, we computed the Lanczos eigenspectrum of the LM loss Hessian restricted
to router parameters (150 iterations per layer, float32).

\begin{table}[h]
\centering
\caption{Hessian eigenspectrum at the uniform-routing solution. Nearly half of
all sampled directions point downhill: uniform routing is a saddle point.}
\label{tab:hessian}
\begin{tabular}{@{}lrrr@{}}
\toprule
Layer & Min $\lambda$ & Max $\lambda$ & Negative fraction \\
\midrule
18 & $-0.050$ & $+0.131$ & 47\% \\
22 & $-0.062$ & $+0.084$ & 48\% \\
27 & $-0.298$ & $+0.158$ & 49\% \\
\bottomrule
\end{tabular}
\end{table}

The analysis revealed two findings. First, the \textbf{root cause} of uniform
routing: our router uses hard argmax, a non-differentiable operation.
$\partial\mathcal{L}/\partial w_{\text{router}} = 0$ exactly---the router
weights are invisible to the optimizer and cannot learn from the LM loss.
Second, when argmax is replaced with a differentiable softmax surrogate, the
Hessian shows abundant negative eigenvalues (47--49\% of sampled directions),
confirming that the LM loss \emph{does} encode routing selectivity. Replacing argmax with
Gumbel-softmax~\cite{jang2017gumbel} ($\tau$: 1.0$\to$0.1) confirms this:
routing selectivity reaches 0.41 at layer~27 (vs.\ $<$0.01 with argmax),
with 82\% of tokens routed with $>$0.3 confidence---all without PPL
degradation (20.14, identical to argmax). This achieves level-2
differentiation (routing) for the first time. However, ablation testing shows
$M_{ij}$ remains uniform (CV~0.12): the routing selectivity does not translate
to causal locality. Instead, the Hessian escape directions lead to
\emph{magnitude asymmetry}---one expert becomes 1.7$\times$ more impactful
than the other across all domains, rather than each expert owning a domain.
Level-2 differentiation (selective routing) is achievable through
differentiable routing; level-3 (causal locality) remains open. A
domain-conditional experiment (4 experts with gradient masking per domain,
supervised router) also failed: the general-domain expert absorbed
$>$99\% of useful learning while domain-specific experts contributed
$<$0.03 PPL each, due to severe data imbalance in C4 ($>$80\% general text).
Even forced specialization on keyword-based domain labels cannot overcome
C4's data imbalance. A \textbf{teacher-driven niche curriculum}
using token-level ZPD scores (Section~\ref{sec:zpd}) achieves strong
\emph{routing} selectivity (92\% of high-ZPD tokens route to the niche
expert), but hot-loading validation shows the expert's contribution is
negligible: removing it causes zero PPL change. The router learned
\emph{where} to send niche tokens, but the rank-4 expert did not develop
enough capacity to become indispensable. This confirms that routing
selectivity (level~2) and causal locality (level~3) are genuinely
distinct: selective routing is achievable through curriculum design, but
selective indispensability requires sufficient expert capacity or longer
training. A top-2 routing variant
with 4 experts (Gumbel-softmax, each token activates 2 of 4) developed a
\emph{capacity hierarchy}---two ``heavy'' experts contributing $\sim$0.56 PPL
each and two ``light'' experts contributing $\sim$0.28---but no domain
selectivity (per-expert CV 0.11--0.16). This is the Standing Committee
pattern~\cite{standing2026} at 4-expert scale: functional importance
differentiates, domain ownership does not. Finally, a damage-prediction
surrogate that directly optimizes for selective indispensability
($L_\text{loc} \propto \text{overlap} - \lambda\,\text{exclusivity}$
of per-token ablation damage) achieves high token-level exclusivity
(1.27 at training time) but the resulting $M_{ij}$ is again
domain-uniform (CV~0.12): expert~1 carries 2.5$\times$ more load than
expert~0, uniformly across all domains. The experts own different
\emph{tokens}, but those tokens are distributed evenly across domains.

\subsection{LoRA Magnitude Analysis}
\label{sec:magnitude}

To understand \emph{why} hot-loading fails, we measured the magnitude of LoRA
deltas relative to the trunk FFN output across 1000 evaluation tokens.

\begin{table}[h]
\centering
\caption{LoRA magnitude analysis at layers 14--27 (Phase 3 checkpoint,
rank-32). Pre-registered prediction: ratio $<$0.05 (small correction).
Result: ratio $\approx$0.25 (substantial correction).}
\label{tab:magnitude}
\begin{tabular}{@{}lrrr@{}}
\toprule
Metric & Expected & Measured & Verdict \\
\midrule
$\|\text{LoRA}\| / \|\text{FFN}\|$ (mean) & $<$0.05 & 0.249 & FAIL \\
Gate CosSim between experts & --- & 0.019 & Orthogonal \\
Routing selectivity (pref.\ vs non-pref.) & $>$2$\times$ & $<$5\% diff & FAIL \\
\bottomrule
\end{tabular}
\end{table}

The LoRA adapters contribute $\sim$25\% of the FFN output magnitude---they
are not small corrections but substantial transformations. Expert weight
matrices are nearly orthogonal (CosSim 0.019), confirming that contrastive
loss produces genuine weight-space differentiation. However, routing
selectivity is weak: both experts produce similar-magnitude deltas regardless
of which tokens are routed to them. This reveals \textbf{redundant representation}: the experts learned different
directions in representation space, but both directions are needed by most
tokens.

We initially hypothesized that this was caused by placing experts in a
domain-agnostic zone (layers 14--17), motivated by the layer divergence
analysis (Section~\ref{sec:divergence}). However, repeating the routing
analysis on the trunk-18 configuration (experts on layers 18--27, where CKA
$>$0.10) yielded the same result: both experts route 49--51\% of tokens
uniformly across all five domains, with per-expert PPL ablation varying by
only $\sim$0.26 points across domains. \textbf{Moving the fork to the
domain-divergent zone did not produce domain-specific routing.}

This suggests that the redundancy is not caused by fork location but by a more
fundamental limitation: contrastive loss differentiates expert \emph{weights}
but does not create functional \emph{routing specialization}. The router
learns to split tokens approximately 50/50 regardless of domain content,
and both experts contribute equally to all domains. Weight-space orthogonality
(CosSim $<$0.3) is a necessary but not sufficient condition for modularity.

% ============================================================================
\section{Related Work}
\label{sec:related}

\paragraph{LoRA and adapters.}
Hu et al.~\cite{hu2022lora} introduced Low-Rank Adaptation for efficient
fine-tuning. We extend LoRA from a fine-tuning tool to a \emph{structural unit
of specialization} within a single model. Li et al.~\cite{li2018intrinsic}
and Aghajanyan et al.~\cite{aghajanyan2021intrinsic} established that
fine-tuning operates on a low intrinsic dimension---our variable-rank mechanism
exploits this directly. HydraLoRA~\cite{hydralora2024} demonstrates that asymmetric residual
specialization (shared A + routed B matrices) is viable, though it does not
test causal modularity (hot-pluggability). MoLoRA~\cite{molora2026} validates
per-token LoRA routing and composability independently.

\paragraph{Shared vs.\ routed experts.}
DeepSeekMoE~\cite{deepseek2024moe} introduced \emph{shared expert isolation}:
always-active experts absorb cross-domain knowledge, freeing routed experts to
specialize. DeepSeek-V3~\cite{deepseekv3} scales this to 1 shared + 256
routed experts. The ``Basic-Refinement Collaboration''
framework~\cite{basicrefinement2025} empirically confirms that shared experts
perform entity recognition and syntactic parsing while routed experts refine
domain-specific associations. Our finding that experts differentiate along
functional axes but not domain axes (Section~\ref{sec:magnitude}) is
independently reported by OpenMoE~\cite{openmoe2024} (context-independent
token-type routing), the Mixtral analysis~\cite{mixtral2024} (``no obvious
patterns based on topic''), and the Standing
Committee~\cite{standing2026} paper (functional experts as stable core, domain
experts as peripheral). The ``Domain vs.\ Driver Experts''
analysis~\cite{domaindriver2026} formalizes this: driver experts handle
function words and syntax, domain experts handle content vocabulary.

\paragraph{Specialization theory.}
B\'{e}na \& Goodman~\cite{bena2025modularity} prove resource constraints are
necessary for specialization---without cost pressure, modules converge. Our
contrastive loss provides one form of constraint; the rank measures the result.
However, our experiments show that weight-space constraints (contrastive loss)
produce representational but not functional specialization. The literature
suggests that top-$K$ routing ($K > 1$) may be required: with $K=1$, each
token must choose a single axis of specialization (typically functional), while
$K \geq 2$ allows simultaneous functional and domain expert
activation~\cite{standing2026, domaindriver2026}.

\paragraph{Low-rank routing.}
L2R~\cite{l2r2026} projects to $r{=}2$ dimensions for routing, confirming that
routing decisions are intrinsically low-dimensional. Our variable-rank adapters
extend this: the routing is low-rank, and so is the specialization itself.

% ============================================================================
\section{Analysis}
\label{sec:analysis}

We address two questions about the emergent structure, reformulated from
our original framing in light of the 2-expert results.

\paragraph{Q1: How does rank affect information efficiency?}
We measure the residual compression ratio (RCR = PPL gain / trainable
parameters) across ranks 1--32 at the trunk-18 boundary.
\begin{table}[h]
\centering
\caption{Rank sweep at trunk-18 boundary (5K steps each, layers 18--27).
Gain measured against frozen-trunk baseline (PPL 26.1, no adapters).
RCR = gain / trainable parameters in millions.}
\label{tab:rank-sweep}
\begin{tabular}{@{}rrrrr@{}}
\toprule
Rank & Params (M) & PPL & Gain & RCR \\
\midrule
1  & 0.164 & 20.73 & 5.42 & \textbf{33.1} \\
2  & 0.328 & 20.65 & 5.49 & 16.8 \\
4  & 0.655 & 20.61 & 5.54 & 8.4 \\
8  & 1.311 & 20.55 & 5.60 & 4.3 \\
16 & 2.621 & 20.43 & 5.71 & 2.2 \\
32 & 5.243 & 20.29 & 5.85 & 1.1 \\
\bottomrule
\end{tabular}
\end{table}

RCR drops 30$\times$ from rank-1 to rank-32, showing that the first few
dimensions of the update-subspace capture most of the learnable signal. In
absolute terms, even rank-32 adds only 5.2M parameters (0.26\% of the 2B base
model) for a 0.44 PPL improvement over rank-1---still highly efficient
relative to the base model. The rank sweep thus validates both the lens and
crystal regimes: low-rank adapters are maximally information-efficient per
parameter, while high-rank adapters remain negligible in total model storage
and can be justified when the marginal quality matters.

\paragraph{Q2: Is the emergent specialization interpretable?}
The first fork produces a structure/content split: Expert~0 routes 71\% of
rare words, Expert~1 routes 86\% of punctuation. This is functional, not
domain-based. The split aligns with the ``Standing Committee''
pattern~\cite{standing2026}: functional experts form a stable core, domain
experts are peripheral. At Qwen3-1.7B scale, domain-level interpretability
has not emerged.

% ============================================================================
\section{Limitations}
\label{sec:limitations}

We identify four limitations that bound our current claims:

\begin{enumerate}
    \item \textbf{Weight differentiation does not imply functional
          specialization.} This is our most important negative finding.
          Contrastive loss reliably produces orthogonal expert weights
          (CosSim $<$0.3) at both fork boundaries tested (layers 14 and 18).
          However, token routing analysis shows both experts route all domains
          uniformly (49--51\% split), and removing either expert degrades all
          domains equally. The experts learn different \emph{directions} in
          representation space, but both directions are needed by all tokens.
          Achieving causal modularity---where expert removal causes
          \emph{selective} domain damage (near-diagonal $M_{ij}$)---requires
          changes to the optimization dynamics. Hessian analysis
          (Section~\ref{sec:hessian}) shows the LM loss \emph{does} encode
          selective routing directions, but auxiliary losses and learning rate
          settings prevent exploitation.
    \item \textbf{No second fork triggered on generic data.} Despite testing
          bimodality thresholds from 0.3 to 0.555, no configuration produced
          more than 2 experts. On generic C4 data, 2-expert systems appear to
          be stable equilibria. Domain-specific curricula may be required to
          trigger deeper forking.
    \item \textbf{Framework tested on one base model only.} All experiments use
          Qwen3-1.7B. Generalization to other architectures (Llama, Mistral,
          GPT-scale) and other data distributions is assumed but not
          demonstrated.
    \item \textbf{Learntropy definition is provisional.} Our use of raw
          cross-entropy as a splitting signal does not distinguish productively
          difficult tokens from noise (Section~\ref{sec:method}). A principled
          learntropy function that models net learning value---closer to
          Wozniak's~\cite{wozniak2018learntropy} definition---may be needed
          for effective curriculum-driven forking.
\end{enumerate}

% ============================================================================
\section{Conclusion}

We set out to test whether specialization depth, adapter rank, and storage tier
correlate strongly. The results are positive on two of three fronts:

\textbf{1.~Information efficiency (achieved).} Variable-rank LoRA adapters on
FFN layers reduce perplexity from 21.20 to 19.83 with 655K trainable
parameters---adapters of $\sim$45\,KB per layer. Layer divergence analysis
empirically identifies the optimal fork boundary (layer~18), and the
accommodation ratio (Eq.~\ref{eq:accommodation}) provides a principled,
measurable criterion for when to expand capacity.

\textbf{2.~Tiered deployment (architecturally valid).} At rank-4, adapter
loads from NVMe consume $<$0.3\% of per-token compute time
(Table~\ref{tab:load-times}). The storage economics align: small adapters are
always cached, larger ones are prefetchable. This has not been tested at scale
(thousands of community-contributed adapters) but the arithmetic is sound.

\textbf{3.~Causal modularity (open problem).} The ablation matrix $M_{ij}$
(Table~\ref{tab:ablation-matrix}) is uniform, not diagonal: removing either
expert degrades all domains equally. Weight orthogonality (CosSim $<$0.3) is
necessary but not sufficient.

Causal locality is not only a modularity requirement---it is also the key to
\emph{compute and memory efficiency}. With uniform $M_{ij}$, every expert must
be active for every token: no expert can be skipped without quality loss, and
tiered storage is architecturally valid but practically useless. With diagonal
$M_{ij}$, irrelevant experts can remain on NVMe until needed, and only a
fraction of experts consume compute per token. Causal locality thus connects
all three pillars: it is what makes compact adapters (pillar~1) and fast
loading (pillar~2) practically beneficial at scale.

Recent literature clarifies that the bottleneck is likely \emph{expert count},
not model scale: domain specialization does not emerge from standard MoE
training at any tested scale up to 34B~\cite{openmoe2024, standing2026}, but
extreme expert counts (262K per layer in Monet~\cite{monet2025}) or explicit
input-space partitioning~\cite{idamoe2025} achieve it at smaller scale. Our
2--4 expert configurations are likely too coarse for domain-level ownership.
A 16-expert rank-1 variant (iso-parameter with 2$\times$rank-4) achieves the
highest $M_{ij}$ CV observed (0.16 vs.\ 0.11--0.14 for 2--4 experts) and
routing selectivity of 0.84, with layer-dependent expert structure---a step
toward but not yet achieving domain selectivity. An initial power law fit ($\text{CV} \approx 0.105 \cdot N^{0.16}$)
predicted $\sim$50 experts for CV$=$0.2, but a 48-expert variant produced
CV$=$0.162---identical to 16 experts. The M$_{ij}$ CV plateaus at $\sim$0.16 for Qwen3-1.7B across all
configurations tested (2--48 experts). A layer divergence analysis on
Qwen3-8B shows 4--5$\times$ higher domain-discriminative fisher ratios
(0.081 vs.\ 0.018), and a preliminary LoRA forking experiment on 8B
achieves $M_{ij}$ CV of 0.52--0.63---\textbf{the first domain-selective
causal locality observed}, with per-expert ablation damage varying
3.5--13.5$\times$ across domains.
The Hessian analysis (Section~\ref{sec:hessian}) confirms the LM loss
\emph{does} encode selective routing directions; the path forward is
differentiable routing (Gumbel-softmax~\cite{jang2017gumbel}) combined with
many more, smaller experts.

% ============================================================================
\appendix

\section{Mathematical Details}
\label{app:math}

\subsection{Learntropy Bimodality Test}

Given an expert $e$ processing tokens $\{x_1, \ldots, x_T\}$, we compute the
per-token loss vector $\ell_e = [\ell_e(x_1), \ldots, \ell_e(x_T)]$ and apply
Hartigan's dip test~\cite{hartigan1985dip} for unimodality. The expert is a
split candidate when the dip statistic exceeds a threshold $\delta$:
\begin{equation}
    \text{DIP}(\ell_e) > \delta \implies \text{split expert } e
\end{equation}
We set $\delta = 0.05$ (the conventional significance level for rejecting
unimodality).

\subsection{Rank Growth Criterion}

After a split, each child adapter begins at rank $r = 1$. Let
$W_{\text{parent}}$ be the parent's weight matrix and $\Delta W = BA$ the LoRA
perturbation at current rank $r$. We compute the relative reconstruction error:
\begin{equation}
    \epsilon_r = \frac{\|W_{\text{target}} - W_{\text{parent}} - BA\|_F}{\|W_{\text{target}} - W_{\text{parent}}\|_F}
\end{equation}
where $W_{\text{target}}$ is the optimal weight matrix estimated from training
data. When $\epsilon_r > \epsilon_{\max}$, the rank is doubled: $r \leftarrow 2r$.

\subsection{Combined Training Objective}

The full training loss for each expert combines the standard language modeling
loss with the contrastive term:
\begin{equation}
    \mathcal{L} = \mathcal{L}_{\text{LM}} + \lambda_c \sum_{i < j} \max\bigl(0,\; \cos(\theta_{ij}) - \tau\bigr)
\end{equation}
We use $\lambda_c = 0.1$ and $\tau = 0.85$ in all experiments.

\section{Shannon's Entropy Argument for Zipf-Aligned Routing}
\label{app:shannon}

Shannon's analysis of English~\cite{shannon1948} established that natural
language has a specific entropy structure: highly redundant on average
($\sim$1 bit per character), but with information concentrated in sparse,
high-surprise moments. This is the source of Zipf's law---common tokens carry
little information and appear frequently, while rare tokens carry much
information and appear rarely.

Standard MoE load-balancing ignores this structure: it allocates equal capacity
to every expert regardless of information content. \tok{}'s argument is a
Shannon argument: parameter allocation should reflect information density. When
specialization depth follows Zipf's law, so does adapter rank, and therefore so
does storage footprint. The most frequently needed adapters are the smallest;
the largest adapters are needed least often. This access pattern is precisely
what tiered storage hierarchies exploit---not by coincidence, but because both
language and memory hierarchies are shaped by power-law economics.

At current scale (Qwen3-1.7B, adapters totaling $<$20\,MB), tiered storage is
unnecessary---everything fits in VRAM. Tiered deployment becomes relevant at
three thresholds: (1)~larger cores (70B+) where VRAM is contested,
(2)~thousands of community-contributed experts, and (3)~edge devices with
6--8\,GB total memory.

\section{Tiered Storage Validation}
\label{app:tiered}

We conducted 121 inference experiments on rented GPU instances (NVIDIA RTX PRO
6000 Blackwell, OLMoE-1B-7B) to calibrate a hardware-realistic cache simulator.

\begin{table}[h]
\centering
\caption{Memory tier economics (Netherlands retail, March 2026, incl.\ VAT).}
\label{tab:tiers}
\begin{tabular}{@{}lrrr@{}}
\toprule
Medium & Capacity & Bandwidth & Cost/GB \\
\midrule
GPU VRAM (GDDR7) & 48\,GB & $\sim$1.8\,TB/s & (fixed at mfg.) \\
DDR5-5600 & 256\,GB & $\sim$90\,GB/s & $\sim$\texteuro{}12/GB \\
NVMe Gen5 & 2--4\,TB & $\sim$14\,GB/s & $\sim$\texteuro{}0.09/GB \\
NVMe Gen4 & 2--4\,TB & $\sim$7\,GB/s & $\sim$\texteuro{}0.07/GB \\
\bottomrule
\end{tabular}
\end{table}

\begin{table}[h]
\centering
\caption{Expert load time from Gen5 NVMe at 14\,GB/s, vs.\ per-token compute
time at 134\,tok/s (7.5\,ms/token).}
\label{tab:load-times}
\begin{tabular}{@{}lrrl@{}}
\toprule
Adapter rank & Size & Load time & Critical path? \\
\midrule
Rank-1 & 4\,KB & 0.3\,$\mu$s & No (CPU cache) \\
Rank-16 & 64\,KB & 4.6\,$\mu$s & No ($<$0.1\% of token time) \\
Rank-64 & 256\,KB & 18\,$\mu$s & No ($<$0.3\% of token time) \\
Full expert (OLMoE) & 16\,MB & 1.1\,ms & Marginal (15\% of token time) \\
\bottomrule
\end{tabular}
\end{table}

Even at rank-64, adapter loads consume $<$0.3\% of a token's compute budget.
Tiered storage is not a compromise---it is invisible at \tok{} adapter sizes.

% ============================================================================
\begin{thebibliography}{99}
\bibitem{hu2022lora} E. Hu et al. LoRA: Low-Rank Adaptation of Large Language Models. \textit{ICLR}, 2022.
\bibitem{geva2021transformer} M. Geva, R. Schuster, J. Berant, O. Levy. Transformer Feed-Forward Layers Are Key-Value Memories. \textit{EMNLP}, 2021.
\bibitem{li2018intrinsic} C. Li et al. Measuring the Intrinsic Dimension of Objective Landscapes. \textit{ICLR}, 2018.
\bibitem{aghajanyan2021intrinsic} A. Aghajanyan et al. Intrinsic Dimensionality Explains the Effectiveness of Language Model Fine-Tuning. \textit{ACL}, 2021.
\bibitem{komatsuzaki2023sparse} A. Komatsuzaki et al. Sparse Upcycling. \textit{arXiv:2212.05055}, 2023.
\bibitem{bena2025modularity} G. B\'{e}na, D. Goodman. Dynamics of Specialization Under Resource Constraints. \textit{Nature Communications}, 2025.
\bibitem{l2r2026} L2R: Low-Rank Lipschitz-Controlled Routing. \textit{arXiv:2601.21349}, 2026.
\bibitem{shannon1948} C. E. Shannon. A Mathematical Theory of Communication. \textit{Bell System Technical Journal}, 27(3):379--423, 1948.
\bibitem{hartigan1985dip} J. A. Hartigan, P. M. Hartigan. The Dip Test of Unimodality. \textit{Annals of Statistics}, 13(1):70--84, 1985.
\bibitem{piaget1952} J. Piaget. \textit{The Origins of Intelligence in Children}. International Universities Press, 1952.
\bibitem{vygotsky1978} L. S. Vygotsky. \textit{Mind in Society: The Development of Higher Psychological Processes}. Harvard University Press, 1978.
\bibitem{wozniak2018learntropy} P. A. Wozniak. Learntropy. \textit{SuperMemo Guru}, 2018. \url{https://supermemo.guru/wiki/Learntropy}.
\bibitem{standing2026} The Illusion of Specialization: Standing Committee in Mixture-of-Experts. \textit{arXiv:2601.03425}, 2026.
\bibitem{deepseek2024moe} X. Dai et al. DeepSeekMoE: Towards Ultimate Expert Specialization. \textit{ACL}, 2024.
\bibitem{deepseekv3} DeepSeek-AI. DeepSeek-V3 Technical Report. \textit{arXiv:2412.19437}, 2024.
\bibitem{hydralora2024} T. Tian et al. HydraLoRA: An Asymmetric LoRA Architecture for Efficient Fine-Tuning. \textit{NeurIPS}, 2024.
\bibitem{molora2026} MoLoRA: Composable Specialization via Per-Token Adapter Routing. \textit{arXiv:2603.15965}, 2026.
\bibitem{basicrefinement2025} Decoding Knowledge Attribution in MoE: A Framework of Basic-Refinement Collaboration. \textit{ACL}, 2025.
\bibitem{openmoe2024} F. Xue et al. OpenMoE: An Early Effort on Open Mixture-of-Experts Language Models. \textit{ICML}, 2024.
\bibitem{mixtral2024} A. Jiang et al. Mixtral of Experts. \textit{arXiv:2401.04088}, 2024.
\bibitem{domaindriver2026} What Gets Activated: Uncovering Domain and Driver Experts in MoE Language Models. \textit{arXiv:2601.10159}, 2026.
\bibitem{jang2017gumbel} E. Jang, S. Gu, B. Poole. Categorical Reparameterization with Gumbel-Softmax. \textit{ICLR}, 2017.
\bibitem{monet2025} S. Park et al. Monet: Mixture of Monosemantic Experts for Transformers. \textit{ICLR}, 2025.
\bibitem{idamoe2025} IDA-MoE: Input Domain Aware Mixture-of-Experts. \textit{arXiv:2510.16448}, 2025.
\end{thebibliography}

\end{document}
